% -----------------------------------------------------------
\section{Knowledge, Know}
% -----------------------------------------------------------
	\label{KNOWLEDGE}
	
% % ----------------------
% \subsection{See also}
% % ----------------------

% ----------------------
\subsection{1828 American Dictionary of the English Language} % *1828
% ----------------------
	\label{KNOWLEDGE-1828}
	
	Note that definition \#1 here is taken straight from Locke, who defines knowledge in the \textit{Essay} this way: ``the perception of the connexion and agreement, or disagreement and repugnancy of any of our Ideas'' (4.1.2; see Yolton's \textit{Locke Dictionary}, 109, PDF 59). So the second part of the definition is taken straight from Locke's \textit{Essay} without modification. Here is part of Yolton's discussion in the \textit{Dictionary}:
	
		\begin{quote}
			Each of the terms in this definition is important for understanding the concept of knowledge with which Locke worked. `Perception' tells us that knowledge is dependent upon one's awareness. Knowledge does not exist apart from us, waiting to be discovered. We tend to say books contain knowledge, we may speak of the book of knowledge, but what books contain is information (or misinformation). Once we have understood that information, we have acquired some knowledge. Similarly for acquiring knowledge about the world from our observation and experience of objects and events: it is our acquisition of ideas about the world which leads to knowledge. Just having ideas is not sufficient for knowledge; we must perceive `the connexion and agreement, or disagreement and repugnancy' of our ideas. It is the relations of and between ideas, when we are aware of those relations, that constitute our knowledge, our knowing this or that. In some cases, we need the help of demonstration before we can perceive the relations between ideas, e.g., the relation between (the ideas of) any three angles of a triangle, equality, and two right angles. In this example, Locke leaves out the phrase `the ideas of, speaking directly of perceiving that the `Equality to two right [angles] does necessarily agree to, and is inseparable from the three Angles of a Triangle'. We may wonder what it is that we perceive when we acquire knowledge: do we perceive ideas or figures? (cf. 4.2.2, 3). In other cases, demonstration is not necessary; we perceive immediately and intuitively without help that, e.g., white is not black ( 4.1.2; 4.2.1). Here, ideas are present in his example: we perceive that the ideas of white and black do not agree (109).
		\end{quote}
		
	It makes me wonder what other definitions in the 1828 dictionary are taken from Locke or others writers like him.

	\begin{compactitem}
		\item[1] A clear and certain perception of that which exists, or of truth and fact; the perception of the connection and agreement, or disagreement and repugnancy of our ideas. We can have no \textit{knowledge} of that which does not exist. God has a perfect \textit{knowledge} of all his works. Human \textit{knowledge} is very limited, and is mostly gained by observation and experience.
		\item[2] Learning; illumination of mind. Ignorance is the curse of God, \textit{knowledge} the wing wherewith we fly to heaven. 
		\item[3] Skill; as a \textit{knowledge} of seamanship.
		\item[4] Acquaintance with any fact or person. I have no \textit{knowledge} of the man or thing.
		\item[5] Cognizance; notice. Ruth 2:10.
		\item[6] Information; power of knowing.
		\item[7] Sexual intercourse. But it is usual to prefix carnal; as carnal \textit{knowledge}.
	\end{compactitem}

% % ----------------------
% \subsection{Oxford English Dictionary} % *OED
% % ----------------------
	% \label{KNOWLEDGE-OED}

	% \begin{compactitem}
		% \item[]
	% \end{compactitem}
	
% ----------------------
\subsection{Restoration Lexicon} % *RL *RESTORATION LEXICON
% ----------------------
	\label{KNOWLEDGE-OED}

	\begin{compactitem}
		\item The cognitive faculties of a human being, referred to either before or after the resurrection. \hyperref[2NEPHI-9-13]{2 Nephi 9:13}.
	\end{compactitem}

% ----------------------
\subsection{Scriptural Usage} % *SCRIPTURES
% ----------------------
	\label{KNOWLEDGE-SCRIPTURES}

	% % -----
	% \subsubsection{Old Testament} % *OT
	% % -----
		% \label{KNOWLEDGE-OT}
	
		% % --
		% \paragraph{KJV}
		% % --
			% \label{KNOWLEDGE-OT-KJV}
	
			% \begin{tabular}{ll}
				% Total occurrences in the OT & 
			% \end{tabular}
			
			% \begin{compactdesc}
				% \item[]
			% \end{compactdesc}
			
		% % --
		% \paragraph{Hebrew Bible}
		% % --
			% \label{KNOWLEDGE-HEBREW}
		
			% %
			% \subparagraph{\texorpdfstring{\he{sample word}}{}}
			% %
				% \label{PAGA} % whatever is in step bible without periods
			
				% \begin{compactdesc}
					% \item[HALOT , PDF ] \textbf{}
						% \begin{compactitem}
							% \item[1]
						% \end{compactitem}
					% \item[BDB , PDF ] \textbf{}
						% \begin{compactitem}
							% \item[1]
						% \end{compactitem}
					% \item[DCH PDF ] \textbf{}
						% \begin{compactitem}
							% \item[1]
						% \end{compactitem}
				% \end{compactdesc}
				
				% \begin{compactdesc}
					% \item[Some OT uses] \textbf{}
						% \begin{compactdesc}
							% \item[]
						% \end{compactdesc}
				% \end{compactdesc}
			
		% % --
		% \paragraph{Septuagint}
		% % --
			% \label{KNOWLEDGE-LXX}
		
			% %
			% \subparagraph{\texorpdfstring{\gr{sample word}}{}}
			% %
		
	% -----
	\subsubsection{New Testament} % *NT
	% -----
		\label{KNOWLEDGE-NT}
	
		% % --
		% \paragraph{KJV}
		% % --
			% \label{KNOWLEDGE-NT-KJV}
			
	
			% \begin{tabular}{ll}
				% Total occurrences in the NT & 
			% \end{tabular}
			
			% \begin{compactdesc}
				% \item[]
			% \end{compactdesc}
			
		% --
		\paragraph{Greek New Testament} % *GREEK
		% --
			\label{KNOWLEDGE-GREEK}
		
			%
			\subparagraph{\texorpdfstring{\gr{οἶδα}}{}}
			%
				\label{OIDA}
			
				\begin{compactdesc}
					\item[BDAG 615a] \textbf{}
						\begin{compactitem}
							\item[1] \textbf{to have information about, \textit{know}}
							\item[1.a] with accusative of person \textit{know someone, know about someone}
							\item[1.b] with accusative of thing
							\item[1.c] with accusative of person and participle in place of the predicate...also without a participle
							\item[1.d] followed by accusative and infinitive
							\item[1.e] followed by \gr{ὅτι}
							\item[1.f] with indirect question following
							\item[1.g] followed by a relative
							\item[1.h] followed by \gr{περί τίνος} \textit{know about something}
							\item[1.i] absolute
							\item[2] \textbf{be intimately acquainted with or stand in a close relation to, \textit{know}}
							\item[3] \textbf{to know/understand how, \textit{can, be able}} with infinitive following
								\begin{compactitem}
									\item James 4:17 is listed as an example of this sense. It also lists \hyperref[MATTHEW-7-11]{Matthew 7:11}, \hyperref[LUKE-11-13]{Luke 11:13}, \hyperref[LUKE-12-56]{Luke 12:56}, \hyperref[PHILIPPIANS-4-12]{Philippians 4:12}, \hyperref[1THESSALONIANS-4-4]{1 Thessalonians 4:4}, \hyperref[1TIMOTHY-3-5]{1 Timothy 3:5}, and \hyperref[2PETER-2-9]{2 Peter 2:9}.
									\item The entry notes that \hyperref[MATTHEW-27-65]{Matthew 27:65} is a related sense but is used absolutely.
								\end{compactitem}
							\item[4] \textbf{to grasp the meaning of something, \textit{understand, recognize, come to know, experience}}
							\item[5] \textbf{to remember, \textit{recollect, recall, be aware of}}
							\item[6] \textbf{to recognize merit, \textit{respect, honor}}
								\begin{compactitem}
									\item Under this sense the entry lists \hyperref[1THESSALONIANS-5-12]{1 Thessalonians 5:12} and \hyperref[EPHESIANS-5-5]{Ephesians 5:5}.
								\end{compactitem}
						\end{compactitem}
					\item[CGL 988b, PDF 1028] \textbf{}
						\begin{compactitem}
							\item[1] have seen or been made aware (of something); \textbf{know} (transitive or intransitive)
							\item[2] (specifically) be aware of (an event, circumstance, item of information, from report or one's own observation); \textbf{know, know about}
							\item[3] be aware (of something) intuitively or through divine revelation; \textbf{know}
							\item[4] be aware of (someone or something) as existing (especially under a certain name); \textbf{know of} or \textbf{about}
							\item[5] be familiar with (the nature or character of someone or something) through experience or acquaintance; \textbf{have experience of, know about, know}
							\item[6] recognise (someone or something) as having authority or requiring obligation; \textbf{acknowledge}
							\item[7] recognize (something) as very important or certain (frequently imperative); \textbf{know} or \textbf{understand for sure}
							\item[8] have understanding or insight; (of a person or god, their mind) \textbf{know}
							\item[9] hold in one's mind (thoughts which are typical of one's character or disposition) \textbf{think}
							\item[10] have practical knowledge or skill (through study or practice); \textbf{be skilled in} ---\textit{crafts, songs, prophecies, deeds of war, athletic contests, medicines} (Homer, Alcmaeon); \textbf{be skilled} ---with genitive \textit{in crafts, use of a weapon, a style of fighting, portents, or similar} (Homer); (of hounds) ---\textit{in the hunt}; \textbf{have the skill} or \textbf{knowledge} ---with infinitive \textit{to do or say something}; \textbf{understand} ---\textit{a subject or topic}
						\end{compactitem}
					\item[LSJ 483a, PDF 554] \textit{I see with the mind's eye}, i.e. \textit{I know}, used a present: pluperfect, I \textit{knew}, used as imperfect
						\begin{compactitem}
							\item[I.1] \textit{know, have knowledge of, be acquainted with}
							\item[I.2] with infinitive, \textit{know how} to do...also, \textit{to be in a condition, be able, have the power}
							\item[I.3] with participle, \textit{to know that} such and such \textit{is} the fact, the participle being in nominative when it is a predicate of the Subject of the Verb
							\item[I.4] less frequently with accusative and infinitive
							\item[I.5] with accusative followed by [words for ``that'']
							\item[I.6] to express disbelief or doubt
							\item[I.7] in similar ellipses with other Conjunctions
							\item[I.8] sometimes parenthetic
						\end{compactitem}
				\end{compactdesc}
				
				\begin{compactdesc}
					\item[Some NT uses] \textbf{}
						\begin{compactdesc}
							\item[]
						\end{compactdesc}
				\end{compactdesc}
				
			%
			\subparagraph{\texorpdfstring{\gr{ἐπιγινώσκω}}{}}
			%
				\label{EPIGINOSKO}
			
				\begin{compactdesc}
					\item[BDAG 325b] \textbf{}
						\begin{compactitem}
							\item[1] \textbf{to have knowledge of something or someone, \textit{know}}
							\item[1a] with the preposition making its influence felt, \textit{know exactly, completely, through and through}
							\item[1b] with no emphasis on the preposition, essentially = \gr{γινώσκω}
							\item[2] \textbf{to ascertain or gain information about something}, with no emphasis on the preposition
							\item[2a] \textbf{\textit{learn, find out}}
							\item[2b] \textbf{\textit{learn to know}}
							\item[2c] \textbf{\textit{notice, perceive, learn of, ascertain}}
							\item[3] \textbf{to connect present information or awareness with what was known before, \textit{acknowledge acquaintance with, recognize, know again}}
							\item[4] \textbf{to indicate that one values the person of another, \textit{acknowledge, give recognition to}}
							\item[5] \textbf{to come to an understanding of, \textit{understand, know}}
						\end{compactitem}
					\item[CGL 549b, PDF 575] \textbf{}
						\begin{compactitem}
							\item[1] understand the identity of, \textbf{recognise}
							\item[2] understand the nature or implications of, \textbf{understand, recognise}
							\item[3] have knowledge of, \textbf{be familiar with, understand}...\textbf{know of}
							\item[4] \textbf{observe, watch, see}...\textbf{detect, find}
							\item[5] \textbf{find out, discover, detect, recognise}
							\item[6] come to a judgement on, \textbf{decide on}...\textbf{decide further}
						\end{compactitem}
					\item[LSJ , PDF ] \textbf{}
						\begin{compactitem}
							\item 
						\end{compactitem}
				\end{compactdesc}
				
				\begin{compactdesc}
					\item[Some NT uses] \textbf{}
						\begin{compactdesc}
							\item[]
						\end{compactdesc}
				\end{compactdesc}	
		
	% % -----
	% \subsubsection{Book of Mormon} % *BOM
	% % -----
		% \label{KNOWLEDGE-BOM}
	
		% \begin{tabular}{ll}
			% Total occurrences in the BOM & 
		% \end{tabular}
		
		% \begin{compactdesc}
			% \item[]
		% \end{compactdesc}
		
	% % -----
	% \subsubsection{Doctrine and Covenants} % *DC
	% % -----
		% \label{KNOWLEDGE-DC}
	
		% \begin{tabular}{ll}
			% Total occurrences in the DC & 
		% \end{tabular}
		
		% \begin{compactdesc}
			% \item[]
		% \end{compactdesc}
		
	% % -----
	% \subsubsection{Pearl of Great Price} % *PGP
	% % -----
		% \label{KNOWLEDGE-PGP}
	
		% \begin{tabular}{ll}
			% Total occurrences in the PGP & 
		% \end{tabular}
		
		% \begin{compactdesc}
			% \item[]
		% \end{compactdesc}
		
% ----------------------
\subsection{Key Phrases} % *KEYPHRASES *PHRASES
% ----------------------
	\label{KNOWLEDGE-PHRASES}

	\begin{compactdesc}
		\item[perfect knowledge] \textbf{20} | Acts 24:22; 2 Nephi 9:13 (2), 14 (2); Jacob 4:12; 7:4; Alma 32:21, 26 (3), 29, 34 (2), 35; 55:1; Ether 3:20; Moroni 7:15, 16, 17
			\begin{compactdesc}
				\item[Bible anchor verses] \phantom{.}
					\begin{compactdesc}
						\item[{\hyperref[ACTS-24-22]{Acts 24:22}}] \gr{Ἀνεβάλετο δὲ αὐτοὺς ὁ Φῆλιξ, ἀκριβέστερον εἰδὼς τὰ περὶ τῆς ὁδοῦ, εἴπας· Ὅταν Λυσίας ὁ χιλίαρχος καταβῇ διαγνώσομαι τὰ καθ᾽ ὑμᾶς·}
					\end{compactdesc}
				\item[Commentary] I don't think the Acts usage of ``perfect knowledge'' is very helpful.
			\end{compactdesc}
	\end{compactdesc}
	
% % ----------------------
% \subsection{Bible Dictionaries} % *DICTIONARIES
% % ----------------------
	% \label{KNOWLEDGE-BD}

% ----------------------
\subsection{Restoration Usage} % *RESTORATION
% ----------------------
	\label{KNOWLEDGE-RESTORATION}

	% -----
	\subsubsection{Joseph Smith} % *JS
	% -----
		\label{KNOWLEDGE-JS}
	
		\begin{compactdesc}
			\item[{\hyperlink{EMS-1843-JS-14-MAY-1843-A-WW-a}{14 May 1843 WW-a}}] The principl of knowledge is the principle of Salvation the Principle can be comprehended, for any one that cannot get knowledge to be Saved will be damned. The Principl of Salvation is given to us through the knowledge of Jesus Christ Salvation is nothing more or less than to triumph over all our enemies \& put them under our feet \& when we have power to put all enemies under our feet in this world \& a knowledge to triumph over all evil spirits in the world to come then we are saved as in the case of Jesus he was to reign untill he had put all enemies under his feet \& the last enemy was death
		\end{compactdesc}
	
	% % -----
	% \subsubsection{Brigham Young} % *BY
	% % -----
		% \label{KNOWLEDGE-BY}
	
		% \begin{compactdesc}
			% \item[DATE/SOURCE]
		% \end{compactdesc}
	
% % ----------------------
% \subsection{Commentary and Studies} % *COMMENTARY *STUDIES
% % ----------------------
	% \label{KNOWLEDGE-COMMENTARY}

	% % -----
	% \subsubsection{Key Points}
	% % -----
		% \label{KNOWLEDGE-KEYS}
	
		% \begin{compactitem}
			% \item
		% \end{compactitem}